\chapter{Introduction}
\label{ch:introduction}

% PRESUPUESTO: ~4 pp (incluida la figura de fallo)
% De lo general a lo específico: robótica -> reconstrucción 3D online -> semántica.

\section{Context}
\label{sec:intro_context}
% Mapeo semántico 3D open-vocabulary: qué es, para qué sirve (robótica, consulta por
% lenguaje natural). Por qué el "online" importa y qué restricciones impone.

\section{Motivation: three failure modes}
\label{sec:intro_motivation}
% Los tres modos de fallo de la propuesta oficial:
%   (1) over-merge       -> instancias distintas de la misma clase, fusionadas
%   (2) fragmentación    -> una instancia vista desde varios puntos de vista, partida
%   (3) objeto compuesto -> partes de un objeto semánticamente unitario, separadas
% FIGURA CLAVE (Rerun): sofá partido en cojines, o instancia contaminada tras drift.
% Esta figura vende la tesis entera; conseguirla pronto.

\section{Problem statement}
\label{sec:intro_problem}
% La tesis en una frase:
%   la decisión de fusión es UNA respuesta binaria a VARIAS preguntas distintas
%   (¿geometría desactualizada? ¿segmentación mala? ¿descriptor ambiguo? ¿uno o dos
%   objetos?), y se toma sobre geometría que el SLAM ya movió por debajo.
% Corolario: un solo umbral no puede arbitrar cuatro preguntas -> por eso cada
% heurística nueva arregla un modo de fallo y rompe otro.

\section{Contributions}
\label{sec:intro_contributions}
% 1. Framework/benchmark controlado de fusión bajo deformación del mapa -> \ref{ch:benchmark}
% 2. Estudio del instance tracking: oportunidades, calidad de máscara, SAM3 -> \ref{ch:instance_tracking}
% 3. Contest: arbitraje de instancias a nivel de punto           -> \ref{ch:contest}
% 4. Consistencia geométrica: refresh por BA local + submapas    -> \ref{ch:geometric_consistency}
% 5. Fusión revisitada: análisis del merge + mejoras del criterio -> \ref{ch:fusion}

\section{Outline}
\label{sec:intro_outline}
% Estructura de la memoria. Explicitar que el orden es CAUSAL, no cronológico:
%   no se puede afinar el criterio si la granularidad es incorrecta;
%   no se puede fijar la granularidad si el sustrato geométrico miente;
%   y no se puede saber nada de nada si no se puede medir una decisión.
